\setquestion
---
tipo: Discursiva
dificuldade: Difícil
disciplina: Química
assunto: pH e pOH
tags: [aquário, amônia, oxigênio dissolvido]
fonte: concurso
concurso: UNICAMP
banca: COMVEST
ano: 2009
numero: "1"
assunto1: pH e pOH
tags1: [pH, amônia, basicidade]
assunto2: pH e pOH
tags2: [oxigênio dissolvido, ppm, cálculo de mols]
resposta1: |
  Aumentar. A amônia atua como base de Brønsted, recebendo H+ da água e liberando íons OH-,
  tornando o meio alcalino. Reação: NH3 + H2O -> NH4+ + OH-.
resposta2: |
  5 ppm em 250 kg de água: m = 5 x (250/1000) = 1,25 g de O2.
  Número de mols: n = m/M = 1,25/32 aprox. 0,039 mol de O2.
---
O aquarismo é uma atividade que envolve a criação de espécies aquáticas em ambiente confinado. O bom funcionamento do aquário depende do controle de uma série de parâmetros, como temperatura, matéria orgânica dissolvida, oxigênio dissolvido, pH, entre outros. Para testar seus conhecimentos químicos, responda às seguintes questões:

Dados: 1 ppm significa que há 1 grama de gás oxigênio dissolvido em 1000 quilogramas de água; massa molar do gás oxigênio = \(32\ \text{g mol}^{-1}\).

\questionitem
Um dos principais produtos do metabolismo dos peixes é a amônia, que é excretada na água. Desconsiderando-se qualquer mecanismo de regulação externa e considerando-se apenas essa excreção de amônia, o valor do pH da água do aquário tende, com o passar do tempo, a aumentar, diminuir ou permanecer constante? Justifique.

\questionitem
Para peixes de água fria, a concentração ideal de gás oxigênio dissolvido na água é de 5 ppm. Considerando-se esse valor e um aquário contendo 250 kg de água, quantos mols de gás oxigênio estão dissolvidos nessa água?

\setquestion
---
tipo: Discursiva
dificuldade: Difícil
disciplina: Química
assunto: Oxirredução
tags: [oxirredução, íons, reação redox]
fonte: concurso
concurso: UNICAMP
banca: COMVEST
ano: 2009
numero: "2"
assunto1: Oxirredução
tags1: [oxirredução, íons, colorimetria]
assunto2: Oxirredução
tags2: [equação química, balanceamento, redox]
resposta1: |
  Com oxigênio: solução escura (castanho), pois I2 é formado ao final das duas etapas.
  Sem oxigênio: solução clara, pois Mn2+ (rosa) e I- (incolor) não reagem.
resposta2: |
  Equação balanceada: Mn4+ + 2I- → Mn2+ + I2.
  (Mn4+ é reduzido a Mn2+ e I- é oxidado a I2.)
---
A quantidade de gás oxigênio dissolvido na água pode ser monitorada através de um teste químico, em que, inicialmente, faz-se o seguinte: a uma amostra de 5 mL de água do aquário, adicionam-se 2 gotas de solução de \ce{Mn^{2+}} e 2 gotas de uma solução de \ce{I^{-}} (em meio básico), agitando-se a mistura. Na seqüência, adiciona-se uma solução para tornar o meio ácido e agita-se a mistura resultante. Sabe-se que em meio básico, o íon \ce{Mn^{2+}} se transforma em \ce{Mn^{4+}} ao reagir com o oxigênio dissolvido na água. Em meio ácido, o \ce{Mn^{4+}} da reação anterior reage com o \ce{I^{-}}, produzindo \ce{I2} e \ce{Mn^{2+}}. Quando não há oxigênio dissolvido, as reações anteriormente descritas não ocorrem.

Dados: \ce{Mn^{2+}} (rosa claro), \ce{Mn^{4+}} (preto), \ce{I^{-}} (incolor), \ce{I2} (castanho escuro).

\questionitem
Correlacione a presença ou a ausência de oxigênio dissolvido com a coloração (clara/escura) do teste. Justifique sua resposta, indicando a espécie responsável pela coloração em cada caso.

\questionitem
Escreva a equação química balanceada para a reação do \ce{Mn^{4+}} com o \ce{I^{-}}, conforme se descreve no texto da questão.

\setquestion
---
tipo: Discursiva
dificuldade: Média
disciplina: Geografia
assunto: Domínios Morfoclimáticos
tags: [Pantanal, pecuária, biodiversidade]
fonte: concurso
concurso: UNICAMP
banca: COMVEST
ano: 2009
numero: "3"
assunto1: Pecuária
tags1: [pecuária intensiva, extensiva]
assunto2: Domínios Morfoclimáticos
tags2: [Pantanal, pastagem, hidrografia]
resposta1: |
  Pecuária intensiva: criação em estábulos, alto investimento em alimentação e medicina veterinária,
  alta produtividade e qualidade do produto final.
  Pecuária extensiva: gado solto em pastagens naturais, baixo investimento, menor produtividade e qualidade.
resposta2: |
  Relevo plano recoberto de gramíneas, solo refertilizado pelas cheias sazonais de verão (clima tropical),
  amplas pastagens naturais e extensa rede hidrográfica (Bacia do Paraguai).
---
O Pantanal já teve 17\% de sua paisagem natural devastados, mas o drama da planície alagada, assim como o de outras áreas úmidas do Brasil, é praticamente ignorado pelos governos estadual e federal, afirmaram cientistas reunidos em Cuiabá para discutir o futuro dessas regiões. Segundo Walfrido Tomás, especialista em gestão de biodiversidade da Embrapa Pantanal, a pecuária intensiva está se difundindo no Pantanal e tem desmatado muito mais do que a tradicional pecuária pantaneira.

\questionitem
Compare as formas de pecuária intensiva e extensiva.

\questionitem
Considerando o Domínio Morfoclimático do Pantanal, quais as características naturais que favorecem a atividade pecuária nessa área?

\setquestion
---
tipo: Discursiva
dificuldade: Média
disciplina: Geografia
assunto: Escala Cartográfica
tags: [abelha africana, escala cartográfica, apicultura]
fonte: concurso
concurso: UNICAMP
banca: COMVEST
ano: 2009
numero: "4"
assunto1: Escala Cartográfica
tags1: [escala, distância, abelha africana]
assunto2: Agricultura
tags2: [apicultura, agricultura familiar, economia]
resposta1: |
  Escala 1:150.000.000 → 1 mm = 150 km. Distância no mapa: 16 mm → 2400 km.
  Não avançou para Patagônia e Andes pelas baixas temperaturas (latitude e altitude),
  inadequadas para a abelha africana adaptada ao clima tropical.
resposta2: |
  Dois aspectos econômicos: (1) complemento de renda com a produção e venda de mel;
  (2) geração de empregos e maior dinamismo do mercado local.
---
A abelha, no Brasil, é um híbrido das abelhas européias (\textit{Apis mellifera mellifera}, \textit{Apis mellifera ligustica}, \textit{Apis mellifera caucasica} e \textit{Apis mellifera carnica}) com a abelha africana (\textit{Apis mellifera scutellata}). Essa abelha, africanizada, possui um comportamento muito semelhante ao da \textit{Apis mellifera scutellata}, em razão da maior adaptabilidade desta raça às condições climáticas do País. Muito agressiva, porém menos que a africana, a abelha do Brasil tem grande facilidade de enxamear, alta produtividade e tolerância a doenças.

\includegraphics{https://www.vestibulares.com.br/Content/data/imagens/78AEE8F24578794598704B3AF1A925D1.png}
\credits{Goudie, A. \textit{The Human Impact on the Natural Environment}. 6ª ed., Malden: Blackwell Publishing, 2006, p. 69.}

\questionitem
Calcule a distância, em quilômetros, de propagação da abelha africana entre o ponto de origem e a cidade de Fortaleza. Por que a propagação da abelha africana não avançou para a Patagônia Argentina e a Cordilheira dos Andes?

\questionitem
A apicultura é uma atividade capaz de causar impactos positivos, tanto sociais quanto econômicos, além de contribuir para a manutenção e preservação de ecossistemas existentes. Aponte dois aspectos econômicos positivos trazidos pela apicultura, em especial para a agricultura familiar.

\setquestion
---
tipo: Discursiva
dificuldade: Difícil
disciplina: Física
assunto: Trabalho e Potência
tags: [trabalho, potência, cavalo-vapor]
fonte: concurso
concurso: UNICAMP
banca: COMVEST
ano: 2009
numero: "5"
assunto1: Trabalho
tags1: [trabalho, gráfico FxD, força constante]
assunto2: Potência
tags2: [potência, cavalo-vapor, conversão de unidades]
resposta1: |
  Área sob a curva com F = 800 N constante entre d = 10 m e d = 50 m.
  tau = F x delta_d = 800 x 40 = 32000 J.
resposta2: |
  P = tau/delta_t = 444000/40 = 11100 W.
  Em CV: 11100/740 = 15 CV.
---
A tração animal pode ter sido a primeira fonte externa de energia usada pelo homem e representa um aspecto marcante da sua relação com os animais.

\includegraphics{https://www.vestibulares.com.br/Content/data/imagens/828EB75DCAD69977EC6737B214F46C6A.png}

\questionitem
O gráfico ao lado mostra a força de tração exercida por um cavalo como função do deslocamento de uma carroça. O trabalho realizado pela força é dado pela área sob a curva \(F \times d\). Calcule o trabalho realizado pela força de tração do cavalo na região em que ela é constante.

\questionitem
No sistema internacional, a unidade de potência é o watt (W) = 1 J/s. O uso de tração animal era tão difundido no passado que James Watt, aprimorador da máquina a vapor, definiu uma unidade de potência tomando os cavalos como referência. O cavalo-vapor (CV), definido a partir da idéia de Watt, vale aproximadamente 740 W. Suponha que um cavalo, transportando uma pessoa ao longo do dia, realize um trabalho total de 444000 J. Sabendo que o motor de uma moto, operando na potência máxima, executa esse mesmo trabalho em 40 s, calcule a potência máxima do motor da moto em CV.

\setquestion
---
tipo: Discursiva
dificuldade: Difícil
disciplina: Física
assunto: Cinemática Vetorial
tags: [campo magnético, vetores, força de resistência]
fonte: concurso
concurso: UNICAMP
banca: COMVEST
ano: 2009
numero: "6"
assunto1: Cinemática Vetorial
tags1: [vetores, campo magnético, deslocamento]
assunto2: Força de Resistência
tags2: [equilíbrio de forças, força de resistência, velocidade constante]
resposta1: |
  O vetor deslocamento segue a direção do campo total (soma vetorial de B_I e B_T).
  Pelo diagrama: componentes 6 e 8 (unidades do grid).
  Módulo: |d| = sqrt(6^2 + 8^2) = sqrt(100) = 10 m.
resposta2: |
  Com velocidade constante: F_res = F_asas. Logo b·v^2 = 0,72.
  v^2 = 0,72 / (5 x 10^-3) = 144, portanto v = 12 m/s.
---
Os pombos-correio foram usados como mensageiros pelo homem no passado remoto e até mesmo mais recentemente, durante a Segunda Guerra Mundial. Experimentos mostraram que seu mecanismo de orientação envolve vários fatores, entre eles a orientação pelo campo magnético da Terra.

\includegraphics{https://www.vestibulares.com.br/Content/Data/imagens/6BB5861A30E3A634C331580A6B11670E.png}

\questionitem
Num experimento, um ímã fixo na cabeça de um pombo foi usado para criar um campo magnético adicional ao da Terra. A figura acima mostra a direção dos vetores dos campos magnéticos do ímã \(\overset{\rightarrow}{B}_{I}\) e da Terra \(\overset{\rightarrow}{B}_{T}\). O diagrama quadriculado representa o espaço em duas dimensões em que se dá o deslocamento do pombo. Partindo do ponto O, o pombo voa em linha reta na direção e no sentido do campo magnético total e atinge um dos pontos da figura marcados por círculos cheios. Desenhe o vetor deslocamento total do pombo na figura e calcule o seu módulo.

\questionitem
Quando em vôo, o pombo sofre a ação da força de resistência do ar. O módulo da força de resistência do ar depende da velocidade \(v\) do pombo segundo a expressão \(F_{res} = bv^{2}\), onde \(b = 5{,}0 \times 10^{-3}\ \text{kg/m}\). Sabendo que o pombo voa horizontalmente com velocidade constante quando o módulo da componente horizontal da força exercida por suas asas é \(F_{asas} = 0{,}72\ \text{N}\), calcule a velocidade do pombo.

\setquestion
---
tipo: Discursiva
dificuldade: Média
disciplina: Biologia
assunto: Mutações
tags: [mutação gênica, evolução, agentes mutagênicos]
fonte: concurso
concurso: UNICAMP
banca: COMVEST
ano: 2009
numero: "7"
assunto1: Evolução
tags1: [mutação, evolução, variabilidade genética]
assunto2: Mutações
tags2: [células germinativas, raios X, hereditariedade]
resposta1: |
  Mutações gênicas produzem novos alelos, gerando variabilidade genética. A seleção natural atua sobre
  essa variabilidade, favorecendo os alelos mais adaptados, resultando no processo de evolução biológica.
resposta2: |
  Raios X podem induzir mutações em células germinativas (precursoras dos gametas). Essas mutações são
  transmitidas aos descendentes, que podem desenvolver com maior frequência tumores ou doenças genéticas.
---
Os animais podem sofrer mutações gênicas, que são alterações na sequência de bases nitrogenadas do DNA. As mutações podem ser espontâneas, como resultado de funções celulares normais, ou induzidas, pela ação de agentes mutagênicos, como os raios X. As mutações são consideradas importantes fatores evolutivos.

\questionitem
Como as mutações gênicas estão relacionadas com a evolução biológica?

\questionitem
Os especialistas afirmam que se deve evitar a excessiva exposição de crianças e de jovens em fase reprodutiva aos raios X, por seu possível efeito sobre os descendentes. Explique por quê.

\setquestion
---
tipo: Discursiva
dificuldade: Média
disciplina: Biologia
assunto: Zoologia
tags: [camuflagem, mimetismo, anfíbios, répteis]
fonte: concurso
concurso: UNICAMP
banca: COMVEST
ano: 2009
numero: "8"
assunto1: Comportamento Animal
tags1: [camuflagem, mimetismo, adaptação]
assunto2: Reprodução Animal
tags2: [anfíbios, répteis, fecundação]
resposta1: |
  Camuflagem: borboletas com coloração dos troncos (ou louva-a-deus assemelhado a folha seca, ou
  bichos-pau a gravetos). Vantagem: esconder dos predadores, aumentando sobrevivência e reprodução.
resposta2: |
  Duas diferenças sapo (anfíbio) vs. cobra (réptil): (1) fecundação externa vs. interna;
  (2) desenvolvimento indireto com fase larval aquática vs. desenvolvimento direto sem fase larval.
  (Aceita-se também: ovos sem casca vs. com casca amniótica.)
---
Ao estudar os animais de uma mata, pesquisadores encontraram borboletas cuja coloração se confundia com a dos troncos em que pousavam mais frequentemente; louva-a-deus e mariposas que se assemelhavam a folhas secas; e bichos-pau semelhantes a gravetos. Observaram que muitas moscas e mariposas assemelhavam-se morfologicamente a vespas e a abelhas e notaram, ainda, a existência de sapos, cobras e borboletas com coloração intensa, variando entre vermelho, laranja e amarelo.

\questionitem
No relato dos pesquisadores estão descritos alguns exemplos de adaptações por eles caracterizadas como mimetismo e camuflagem. Identifique no texto um exemplo de camuflagem. Explique uma vantagem dessas adaptações para os animais.

\questionitem
No texto são citados vários animais, entre eles sapos e cobras. Esses animais pertencem a grupos de vertebrados que apresentam diferenças relacionadas com a reprodução. Indique duas dessas diferenças.

\setquestion
---
tipo: Discursiva
dificuldade: Difícil
disciplina: Matemática
assunto: Função Afim
tags: [função afim, coeficiente angular, gráfico]
fonte: concurso
concurso: UNICAMP
banca: COMVEST
ano: 2009
numero: "9"
assunto1: Função Afim
tags1: [coeficiente angular, gráfico, função afim]
assunto2: Função Afim
tags2: [função afim, cálculo acumulado, proporcionalidade]
resposta1: |
  Coeficiente angular: m = (34200 - 0)/(87 - 72) = 2280 baleias/ano (soviéticos/russos).
  Gráfico: zero até 1972, reta crescente 1973–1987 (2280/ano), constante em 34200 de 1988 a 2005.
resposta2: |
  Japão: 6840 baleias em 18 anos → proporcionalmente 1140 entre 1987 e 1990.
  Total 1965–1990: Brasil 13500 + Rússia 34200 + Japão (35000 + 1140) = 83840 baleias.
---
Na década de 1960, com a redução do número de baleias de grande porte, como a baleia azul, as baleias minke antárticas passaram a ser o alvo preferencial dos navios baleeiros que navegavam no hemisfério sul. O gráfico ao lado mostra o número acumulado aproximado de baleias minke antárticas capturadas por barcos japoneses, soviéticos/russos e brasileiros, entre o final de 1965 e o final de 2005.

\questionitem
No gráfico abaixo, trace a curva que fornece o número aproximado de baleias caçadas anualmente por barcos soviéticos/russos entre o final de 1965 e o final de 2005. Indique também os valores numéricos associados às letras A e B apresentadas no gráfico, para que seja possível identificar a escala adotada para o eixo vertical.

\includegraphics{https://www.vestibulares.com.br/Content/Data/imagens/509915D32B7BE931250DE9CB8B17A36C.png}

\questionitem
Calcule o número aproximado de baleias caçadas pelo grupo de países indicado no gráfico entre o final de 1965 e o final de 1990.

\includegraphics{https://www.vestibulares.com.br/Content/Data/imagens/734F72B55D8795BFB227277BE30DED7F.png}

\setquestion
---
tipo: Discursiva
dificuldade: Média
disciplina: Matemática
assunto: Sistemas Lineares
tags: [sistema de equações, paralelepípedo, piscicultura]
fonte: concurso
concurso: UNICAMP
banca: COMVEST
ano: 2009
numero: "10"
assunto1: Sistemas Lineares
tags1: [sistema de equações, duas variáveis]
assunto2: Geometria Espacial
tags2: [volume, paralelepípedo, densidade]
resposta1: |
  Sistema: a + b = 600 e 1,5a + b = 800. Subtração: 0,5a = 200 → a = 400 (espécie A), b = 200 (espécie B).
resposta2: |
  V = 2L^2 (largura = comprimento = L, altura = 2 m).
  400 = 7200/(2L^2) → L^2 = 9 → L = 3 m. Dimensões mínimas: 3 m x 3 m x 2 m.
---
Em um sistema de piscicultura superintensiva, uma grande quantidade de peixes é cultivada em tanques-rede colocados em açudes, com alta densidade populacional e alimentação à base de ração. Os tanques-rede têm a forma de um paralelepípedo e são revestidos com uma rede que impede a fuga dos peixes, mas permite a passagem da água.

\questionitem
Um grupo de 600 peixes de duas espécies foi posto em um conjunto de tanques-rede. Os peixes consomem, no total, 800 g de ração por refeição. Sabendo-se que um peixe da espécie A consome 1,5 g de ração por refeição e que um peixe da espécie B consome 1,0 g por refeição, calcule quantos peixes de cada espécie o conjunto de tanques-rede contém.

\questionitem
Para uma determinada espécie, a densidade máxima de um tanque-rede é de 400 peixes adultos por metro cúbico. Suponha que um tanque possua largura igual ao comprimento e altura igual a 2 m. Quais devem ser as dimensões mínimas do tanque para que ele comporte 7200 peixes adultos da espécie considerada?
